\begin{center}[18.]\end{center}
If one adds to \textit{the masses inside a solid region} the surface masses that are equivalent to them in the exterior, with opposite signs, one obtains a body carrying positive and negative masses. Together these exert no attractive force at all on any point in the exterior.

\SourceLeaf{629}
It is easy to prove

1.\quad that this body also remains in equilibrium as a result of the reaction of external masses

2.\quad that if the masses in question are magnetic fluids, the body also exerts no force at all on a rheophore.

It is harder to prove

3.\quad that the body also remains in equilibrium despite the reaction of the rheophore.

The last point rests on the following considerations.

Let $\mathrm dm$ be an element of the body with coordinates $x,y,z$. The symbol $\int$ refers to all $\mathrm dm$. Let $\mathrm ds$ be an element of a galvanic current, with $a,b,c;\ a+\mathrm da,\ b+\mathrm db,\ c+\mathrm dc$ the coordinates of its endpoints. Thus $a,b,c$ are functions of $s$. The current intensity is $=1$. The symbol $S$ denotes summation with respect to $\mathrm ds$.

If $X,Y,Z$ are the components of the total accelerating force acting on $\mathrm dm$, then, as is well known, the conditions for equilibrium are
\[
\begin{gathered}
\int X\mathrm dm,\int Y\mathrm dm,\int Z\mathrm dm,\int(zY-yZ)\mathrm dm,\\
\int(xZ-zX)\mathrm dm,\int(yX-xY)\mathrm dm\text{ all }=0
\end{gathered}
\]

Put $(x-a)^2+(y-b)^2+(z-c)^2=rr$ (positive $x,y,z$ point forward, to the right, and upward, respectively).
\[
\begin{aligned}
X&=S\frac{(y-b)\mathrm dc-(z-c)\mathrm db}{r^3},\\
Y&=S\frac{(z-c)\mathrm da-(x-a)\mathrm dc}{r^3},\\
Z&=S\frac{(x-a)\mathrm db-(y-b)\mathrm da}{r^3}
\end{aligned}
\]
\[
\text{therefore}\qquad\int X\mathrm dm=S(\eta\mathrm dc-\zeta\mathrm db)
\]
\[
\begin{aligned}
&\int(yX-xY)\mathrm dm\\
&=SV\mathrm dc+S\{(\xi c-\zeta a)\mathrm da+(\eta c-\zeta b)\mathrm db+(a\xi+b\eta+c\zeta)\mathrm dc\}\\
&\qquad-\int z\mathrm dmS\frac{\mathrm d\frac1r}{\mathrm ds}\mathrm ds
\end{aligned}
\]
[Here $\displaystyle\int\frac{\mathrm dm}{r}=V,\quad\int\frac{x-a}{r^3}\mathrm dm=\xi$ and so forth. All these quantities vanish.]

\medskip
1836 February 18.

The derivation is much simpler in the following form.

If $x,y,z;\ x+\mathrm dx,\ y+\mathrm dy,\ z+\mathrm dz$ are the coordinates of any two points infinitely close to each other, then the variation of
\[
\mathrm dx^2+\mathrm dy^2+\mathrm dz^2
\]
\SourceLeaf{630}
must, regardless of the values of $\mathrm dx,\mathrm dy,\mathrm dz,$ become $=0$. For brevity we denote $\delta x,\delta y,\delta z$ by $\alpha\xi,\alpha\eta,\alpha\zeta$, where $\alpha$ denotes a constant infinitesimal coefficient. Thus
\[
\begin{aligned}
\delta(\mathrm dx^2+\mathrm dy^2+\mathrm dz^2)&=2\alpha(\mathrm dx.\mathrm d\xi+\mathrm dy.\mathrm d\eta+\mathrm dz.\mathrm d\zeta)\\
&=2\alpha\Big\{\frac{\mathrm d\xi}{\mathrm dx}.\mathrm dx^2+\frac{\mathrm d\eta}{\mathrm dy}.\mathrm dy^2+\frac{\mathrm d\zeta}{\mathrm dz}.\mathrm dz^2\\
&\quad+\left(\frac{\mathrm d\xi}{\mathrm dy}+\frac{\mathrm d\eta}{\mathrm dx}\right)\mathrm dx\mathrm dy+\left(\frac{\mathrm d\xi}{\mathrm dz}+\frac{\mathrm d\zeta}{\mathrm dx}\right)\mathrm dx\mathrm dz\\
&\quad+\left(\frac{\mathrm d\eta}{\mathrm dz}+\frac{\mathrm dz}{\mathrm dy}\right)\mathrm dy\mathrm dz\Big\}
\end{aligned}
\]
Clearly, then, the following must hold.
\[
\begin{aligned}
1.\quad0&=\frac{\mathrm d\xi}{\mathrm dx}\\
2.\quad0&=\frac{\mathrm d\eta}{\mathrm dy}\\
3.\quad0&=\frac{\mathrm d\zeta}{\mathrm dz}\\
4.\quad\frac{\mathrm d\xi}{\mathrm dy}&=-\frac{\mathrm d\eta}{\mathrm dx}=r\\
5.\quad\frac{\mathrm d\zeta}{\mathrm dx}&=-\frac{\mathrm d\xi}{\mathrm dz}=q\\
6.\quad\frac{\mathrm d\eta}{\mathrm dz}&=-\frac{\mathrm d\zeta}{\mathrm dy}=p
\end{aligned}
\]
From (6) (4) (5) (6), from (6) (2), and from (6) (3), one obtains
\[
\begin{aligned}
\frac{\mathrm dp}{\mathrm dx}&=\frac{\mathrm dd\eta}{\mathrm dx\mathrm dz}=-\frac{\mathrm dd\xi}{\mathrm dz\mathrm dy}=\frac{\mathrm dd\zeta}{\mathrm dy\mathrm dx}=-\frac{\mathrm dp}{\mathrm dx}=0\\
\frac{\mathrm dp}{\mathrm dy}&=\frac{\mathrm dd\eta}{\mathrm dy\mathrm dz}=0\\
\frac{\mathrm dp}{\mathrm dz}&=-\frac{\mathrm dd\zeta}{\mathrm dz\mathrm dy}=0
\end{aligned}
\]
so $p$ is constant.

Likewise, it follows that $q,r$ are constant,

and therefore
\[
\begin{aligned}
\xi&=a+ry-qz\\
\eta&=b+pz-rx\\
\zeta&=c+qx-py
\end{aligned}
\]

\begin{center}[19.]\par\medskip\textit{Proof of} \textsc{Ampère’s} \textit{fundamental theorem.}\end{center}
Let positive magnetic fluid be distributed uniformly over a bounded surface (I), with $x,y,z$ the coordinates of an arbitrary point on it, so that the amount of magnetic fluid per unit area is $=k$.

Let an element of a galvanic current of intensity $i$ extend from $0,0,0$ to $0,0,\zeta$.

\SourceLeaf{631}
For brevity, write
\[
r=\BareRoot{(xx+yy+zz)}
\]

An element $\omega$ of the $x,y$ coordinate plane corresponds to the element of surface (I) $\, .\, .\, .\,\omega\BareRoot{\left(1+\left(\dfrac{\mathrm dz}{\mathrm dx}\right)^2+\left(\dfrac{\mathrm dz}{\mathrm dy}\right)^2\right)}$ and to the amount of magnetic fluid $k\omega\BareRoot{\left(1+\left(\dfrac{\mathrm dz}{\mathrm dx}\right)^2+\left(\dfrac{\mathrm dz}{\mathrm dy}\right)^2\right)}$. The action of this fluid on the current element $i\zeta$ thus resolves into three component forces
\[
\begin{gathered}
ik\zeta\omega.\frac{y}{r^3}\BareRoot{\left(1+\left(\frac{\mathrm dz}{\mathrm dx}\right)^2+\left(\frac{\mathrm dz}{\mathrm dy}\right)^2\right)}\\
-ik\zeta\omega.\frac{x}{r^3}\BareRoot{\left(1+\left(\frac{\mathrm dz}{\mathrm dx}\right)^2+\left(\frac{\mathrm dz}{\mathrm dy}\right)^2\right)}\\
0
\end{gathered}
\]

Let us now take a second surface (II) parallel to (I), an infinitesimal distance $\varepsilon$ below it. That is, to each point $x,y,z$ in (I) there corresponds in II the point given by
\[
\begin{gathered}
x+\frac{\varepsilon}{\BareRoot{\left(1+\left(\frac{\mathrm dz}{\mathrm dx}\right)^2+\left(\frac{\mathrm dz}{\mathrm dy}\right)^2\right)}}.\left(\frac{\mathrm dz}{\mathrm dx}\right)\\
y+\frac{\varepsilon}{\BareRoot{\left(1+\left(\frac{\mathrm dz}{\mathrm dx}\right)^2+\left(\frac{\mathrm dz}{\mathrm dy}\right)^2\right)}}.\left(\frac{\mathrm dz}{\mathrm dy}\right)\\
z-\frac{\varepsilon}{\BareRoot{\left(1+\left(\frac{\mathrm dz}{\mathrm dx}\right)^2+\left(\frac{\mathrm dz}{\mathrm dy}\right)^2\right)}}
\end{gathered}
\]

Let negative magnetic fluid be distributed over this second surface in such a way that each part of II contains just as much negative fluid as the corresponding part of I contains positive fluid. The combined effect of the fluids carried by the corresponding elements of I and II is therefore given by
\[
\begin{aligned}
\frac{\varepsilon ik\zeta.\omega}{r^5}\left\{3xy.\left(\frac{\mathrm dz}{\mathrm dx}\right)+(3yy-rr)\left(\frac{\mathrm dz}{\mathrm dy}\right)-3yz\right\}&=X\omega\\
-\frac{\varepsilon ik\zeta.\omega}{r^5}\left\{(3xx-rr)\left(\frac{\mathrm dz}{\mathrm dx}\right)+3xy\left(\frac{\mathrm dz}{\mathrm dy}\right)-3xz\right\}&=Y\omega\\
0&=Z\omega
\end{aligned}
\]

Let us begin by transforming $\omega X$, which we first put in the form
\SourceLeaf{632}
\[
\begin{aligned}
\omega X&=\frac{\varepsilon ik\zeta\omega}{r^5}\left\{3xy\left(\frac{\mathrm dz}{\mathrm dx}\right)+(2rr-3xx-3zz)\left(\frac{\mathrm dz}{\mathrm dy}\right)-3yz\right\}\\
&=\frac{\varepsilon ik\zeta\omega}{r^5}\Big\{\left(x\frac{\mathrm ddz}{\mathrm dx\mathrm dy}+\frac{\mathrm dz}{\mathrm dy}\right)rr-3\left(x+z\frac{\mathrm dz}{\mathrm dx}\right)x\frac{\mathrm dz}{\mathrm dy}\\
&\quad-\left(x\frac{\mathrm ddz}{\mathrm dx\mathrm dy}-\frac{\mathrm dz}{\mathrm dy}\right)rr+3\left(y+z\frac{\mathrm dz}{\mathrm dy}\right)\left(x\frac{\mathrm dz}{\mathrm dx}-z\right)\Big\}\\
&=\varepsilon ik\zeta\omega\left\{\frac{\mathrm d\dfrac{x\left(\frac{\mathrm dz}{\mathrm dy}\right)}{r^3}}{\mathrm dx}-\frac{\mathrm d\dfrac{x\left(\frac{\mathrm dz}{\mathrm dx}\right)-z}{r^3}}{\mathrm dy}\right\}
\end{aligned}
\]

To determine the total effect parallel to the $x$-axis, we call the projection of I onto the $x,y$-plane, or the collection of all $\omega$, III. We then have to find the integral $\int\omega X$ taken over all $\omega$. In other words, if we write $\mathrm dx\mathrm dy$ in place of $\omega$, we have to carry out the double integration of $X\mathrm dx\mathrm dy$.

Here one must consider
\[
\varepsilon ik\zeta\mathrm dy.\left\{\frac{\mathrm d\dfrac{x\left(\frac{\mathrm dz}{\mathrm dy}\right)}{r^3}}{\mathrm dx}.\mathrm dx\right\}
\]
and
\[
\varepsilon ik\zeta\mathrm dx.\left\{\frac{\mathrm d\dfrac{x\left(\frac{\mathrm dz}{\mathrm dx}\right)-z}{r^3}}{\mathrm dy}.\mathrm dy\right\}
\]
separately. For the first expression, one divides III into infinitely many infinitely narrow strips parallel to the $x$-axis. For the second, one divides it into similar strips parallel to the $y$-axis. It then follows very easily that the total is
\[
\varepsilon ik\zeta\int\frac{x\mathrm dz-z\mathrm dx}{r^3}=\varepsilon ik\zeta\int\frac{x\dfrac{\mathrm dz}{\mathrm ds}-z\dfrac{\mathrm dx}{\mathrm ds}}{r^3}.\mathrm ds
\]
with the integration extending over the entire boundary of III, traversed in the direction that keeps III on the right.

In a similar way, one obtains for the sum of all $Y\omega$
\[
\varepsilon ik\zeta\int\frac{y\mathrm dz-z\mathrm dy}{r^3}=\varepsilon ik\zeta\int\frac{y\dfrac{\mathrm dz}{\mathrm ds}-z\dfrac{\mathrm dy}{\mathrm ds}}{r^3}.\mathrm ds
\]
This, together with [No. 14], proves \textsc{Ampère}’s law.

